\begin{frame}{역사: Elman 네트워크와 ``시간 속의 구조''}
  RNN의 원형은 1990년 인지과학자 제프리 엘먼(Jeffrey Elman)의 논문
  ``Finding Structure in Time''(Elman, 1990).
  \begin{itemize}
    \item 질문: 신경망이 단어 목록을 외우는 대신 \textbf{시간에 걸친
      패턴}을 추출할 수 있을까?
    \item 구조: 은닉층을 \kb{센트럴 유닛}(지금 입력 처리)과
      \kb{라테럴 유닛}(이전 상태의 복사본 저장)으로 나눔
    \item 라테럴 유닛의 출력을 \textbf{다음 시점의 입력으로} 다시
      돌린다 --- 지금의 RNN과 거의 똑같음
    \item 단순한 반복 연결을 ``시간을 기억하는 장치''로 재해석한 것이
      핵심
  \end{itemize}
  \vskip0.5em
  \begin{block}{유명한 실험과 계보}
    같은 어미 ``-pits''가 ``pit(구덩이)''의 복수인지 ``pitch(공)''의
    복수인지를 문맥이 구분하는 듯이 행동한다는 관찰 (나중에 비판도
    받았으나 ``같은 토큰도 맥락에 따라 다른 내적 표현'' 직관의 원형).
    1997년 호흐하이터·슈미트허버가 제안한 \kb{LSTM} 이 2010년대
    음성인식·번역에서 비약적 성능을 보여주며 RNN이 비로소 주류가 됨
    (11.3절에서 다시 만남).
  \end{block}
\end{frame}
